\scanpage{8}
% archive foliation: 7 (upper right and lower left)
\[
F|\mathcal{M}^{+0}(S)\simeq F_{\unsure{\xi|\zeta}}
\]
où $F_{\unsure{\xi|\zeta}}$ est le foncteur évident sur $\mathcal{M}^{+0}(S)$ : fibre en $\unsure{\xi|\zeta}$.
(NB. $S$ étant connexe, ces foncteurs sont isom entre eux, donc la \uline{condition} d'association ne dépend pas de $\unsure{\xi|\zeta}$. Mais il faut le considérer plutôt comme une \uline{donnée} supplémentaire, comme \uline{donnée} de l'isom en question).

\uline{Ex 1} Si source $\unsure{\xi|\zeta}=\operatorname{Spec}(\mathbb{C})$, \deleted{alors} $F=F_{\unsure{\xi|\zeta}}^{\mathrm{Betti}}$, alors on a une telle donnée.
\inserted{\uline{Ex. 2}\quad $T^\ell_{\unsure{\xi|\zeta}}$.}

\uline{Ex 3} Si $S=\operatorname{Spec}k$, $k$ corps de car.\ nulle (p.ex.\ $k=\mathbb{Q}$) et si $F$ est le foncteur fibre \deleted{sur}\inserted{à coefficients dans} $k$ $T_{\mathrm{DR}}$, \deleted{alors} ou $T_{\mathrm{Hdg}}$, alors il n'est pas en général associé aux pts géométriques.

\uline{Ex 4} Si \deleted{$k$ est} $S$ est le spectre d'un corps alg clos, alors $\forall$ $F$ est associé à $\unsure{\xi|\zeta}$, du moins si $k$ est dénombrable, \inserted{i.e.\ de degré de transc.\ absolu dénombrable}, donc $\mathcal{M}(k)$ à engendrement dénombrable\dots

\deleted{Alors on trouve un isom canonique
\[
\left\{\begin{aligned}
\pi_1(S,\unsure{\xi|\zeta})&\simeq\varprojlim G_\alpha/G_\alpha^0\\
\text{les }G_\alpha/G_\alpha^0&\text{ sont des groupes constants}\\
&\text{sur }K\ \dots
\end{aligned}\right.
\]
\unsure{Cela|Ceci} résulte du \deleted{fait que} $\otimes$-foncteur canonique
\diagram{\begin{tikzcd}
\operatorname{Mod}\,\unsure{\mathrm{fin}|\mathrm{cont}}(\pi_1)\simeq \mathcal{M}^{+0}(S)
\arrow[r] & \mathcal{M}(S)
\end{tikzcd}}
qui définit
\diagram{\begin{tikzcd}
G\arrow[r]&\pi_1
\end{tikzcd}}
}
